\chapter{Research aims and questions}
\label{chap:research_aims_and_questions}

\section{Primary aim}
\label{sec:primary_aim}

The primary aim was to identify internally reproducible within-patient proteomic differences between tumour and matched non-tumour tissue while preserving the uncertainty created by missing values, high dimensionality and incomplete project metadata.

\section{Secondary aims}
\label{sec:secondary_aims}

The secondary aims were to:

\begin{enumerate}
\item construct an immutable and auditable data model from the source workbook;
\item characterise specimen quality, pair concordance, multivariate structure and missingness before differential testing;
\item estimate paired abundance changes and paired detection differences with multiplicity control;
\item organise feature-level evidence into versioned biological pathways without converting enrichment into causal claims;
\item assess whether stable patient response subgroups are supported;
\item test tissue-state predictability using leakage-controlled machine learning and appropriate baselines;
\item quantify sensitivity to analytical choices and individual patients;
\item integrate inference, pathways and predictive stability without allowing one evidence type to substitute for another; and
\item freeze a prospective validation-readiness package that cannot be redesigned after future outcomes are observed.
\end{enumerate}

\section{Research questions}
\label{sec:research_questions}

The analysis addressed the following questions:

\begin{itemize}
\item Which source features exhibit reproducible quantitative tumour-minus-non-tumour changes?
\item Which features show reproducible paired detection asymmetry?
\item Which Reactome pathways organise these changes across both ranked and patient-level views?
\item Does the patient-change geometry support stable molecular response clusters?
\item Can tumour and matched non-tumour specimens be distinguished using models evaluated on entirely held-out patients?
\item Are feature, pathway and prediction conclusions robust to plausible representations, cohort definitions, complete-pair thresholds, estimators, seeds and patient deletions?
\item Which findings are sufficiently robust for prospective follow-up, and what evidence is still missing?
\end{itemize}
